%%%%%%%%%%%%%%%%%%%%%%%%%%%%%%%%%%%%%%%%%%%%%%%%%%%%%%%%%%%%
%% Multi-Agent Game Theory at Scale
%% Research Notes — ETH Zurich ML Lecture Series
%% Agentic AI Frontier, Lecture 04
%%%%%%%%%%%%%%%%%%%%%%%%%%%%%%%%%%%%%%%%%%%%%%%%%%%%%%%%%%%%
\documentclass[11pt,a4paper]{article}

%%% ─── packages ──────────────────────────────────────────
\usepackage{amsmath,amsthm,amssymb,mathtools}
\usepackage[margin=2.5cm]{geometry}
\usepackage{hyperref}
\usepackage{enumitem}
\usepackage{booktabs}
\usepackage{xcolor}
\usepackage{tikz}
\usepackage{algorithm}
\usepackage{algpseudocode}
\usepackage{natbib}
\usepackage{thmtools}

%%% ─── theorem environments ─────────────────────────────
\theoremstyle{plain}
\newtheorem{theorem}{Theorem}[section]
\newtheorem{lemma}[theorem]{Lemma}
\newtheorem{proposition}[theorem]{Proposition}
\newtheorem{corollary}[theorem]{Corollary}
\theoremstyle{definition}
\newtheorem{definition}[theorem]{Definition}
\newtheorem{example}[theorem]{Example}
\newtheorem{remark}[theorem]{Remark}
\newtheorem{openproblem}{Open Problem}

%%% ─── macros ────────────────────────────────────────────
\newcommand{\R}{\mathbb{R}}
\newcommand{\E}{\mathbb{E}}
\newcommand{\N}{\mathbb{N}}
\newcommand{\cA}{\mathcal{A}}
\newcommand{\cS}{\mathcal{S}}
\newcommand{\cH}{\mathcal{H}}
\newcommand{\cI}{\mathcal{I}}
\newcommand{\cZ}{\mathcal{Z}}
\newcommand{\br}{\mathrm{BR}}
\newcommand{\NE}{\mathrm{NE}}
\newcommand{\supp}{\mathrm{supp}}
\newcommand{\KL}{\mathrm{KL}}
\DeclareMathOperator*{\argmax}{arg\,max}
\DeclareMathOperator*{\argmin}{arg\,min}

%%% ─── metadata ──────────────────────────────────────────
\title{\textbf{Multi-Agent Game Theory at Scale:\\
Mathematical Foundations for Agentic AI}}
\author{Research Notes --- ETH Z\"urich ML Lecture Series\\[4pt]
\small Agentic AI Frontier, Lecture~04}
\date{Spring Semester 2026}

\begin{document}
\maketitle

\begin{abstract}
These notes provide a self-contained treatment of multi-agent game theory
as it applies to modern large-scale AI systems.  We develop the
mathematical foundations of Nash equilibria, counterfactual regret
minimisation (CFR), and mean-field games, then trace the line from
classical solution concepts through landmark systems---Libratus, Pluribus,
Cicero---to contemporary questions about multi-agent LLM deployments and
AI safety.  The target audience is Master's students with background in
machine learning, optimisation, and probability theory.
\end{abstract}

\tableofcontents
\newpage

%===================================================================
\section{Introduction}\label{sec:intro}
%===================================================================

Game theory studies strategic interaction: settings where the outcome for
each participant depends on the actions of all.  As AI systems move from
single-agent pipelines to \emph{agentic} architectures---where multiple
LLM-powered agents negotiate, plan, and act in shared
environments---the classical toolkit of Nash equilibria, regret
minimisation, and mechanism design becomes directly relevant.

\paragraph{Why game theory for ML researchers.}
\begin{enumerate}[label=(\roman*)]
  \item \textbf{Training as a game.}  GANs, RLHF, and self-play all
    involve multiple interacting objectives.
  \item \textbf{Deployment as a game.}  Autonomous agents bidding in
    auctions, negotiating contracts, or coordinating in multi-agent
    environments face genuine strategic uncertainty.
  \item \textbf{Safety as a game.}  Alignment can be formalised as a
    principal--agent problem; robustness against adversarial attacks is a
    zero-sum game.
\end{enumerate}

The notes are organised as follows.  Section~\ref{sec:nash} reviews
classical solution concepts.  Section~\ref{sec:cfr} develops
counterfactual regret minimisation.
Sections~\ref{sec:libratus}--\ref{sec:cicero} cover the landmark poker
and Diplomacy AI systems.  Section~\ref{sec:mfg} introduces mean-field
games.  Section~\ref{sec:opponent_shaping} treats opponent shaping.
Section~\ref{sec:agentic} draws connections to modern agentic AI.
Section~\ref{sec:open} lists open problems, and
Section~\ref{sec:biblio} provides an annotated bibliography.

%===================================================================
\section{Nash Equilibria and Classical Solution Concepts}\label{sec:nash}
%===================================================================

\subsection{Normal-Form Games}

\begin{definition}[Normal-form game]
A \emph{finite normal-form game} is a tuple
$\Gamma = (N, (A_i)_{i\in N}, (u_i)_{i\in N})$ where
\begin{itemize}
  \item $N = \{1,\dots,n\}$ is the set of \emph{players},
  \item $A_i$ is the finite \emph{action set} of player~$i$,
  \item $u_i : \prod_{j} A_j \to \R$ is the \emph{utility function} of
    player~$i$.
\end{itemize}
\end{definition}

A \emph{mixed strategy} for player~$i$ is a distribution
$\sigma_i \in \Delta(A_i)$.  Write $\sigma = (\sigma_i)_{i\in N}$ for a
\emph{strategy profile} and $\sigma_{-i}$ for the profile of all players
except~$i$.  The expected utility under $\sigma$ is
\[
  U_i(\sigma)
  = \sum_{a \in \prod_j A_j}
    \Bigl(\prod_{j\in N} \sigma_j(a_j)\Bigr)\, u_i(a).
\]

\subsection{Nash Equilibrium}

\begin{definition}[Nash equilibrium]
A strategy profile $\sigma^*$ is a \emph{Nash equilibrium} (NE) if, for
every player~$i$ and every alternative strategy $\sigma_i' \in
\Delta(A_i)$,
\[
  U_i(\sigma_i^*, \sigma_{-i}^*) \;\geq\; U_i(\sigma_i', \sigma_{-i}^*).
\]
Equivalently, every player plays a \emph{best response} to the others.
\end{definition}

\begin{definition}[$\varepsilon$-Nash equilibrium]
A profile $\sigma$ is an $\varepsilon$-Nash equilibrium if for every
player~$i$ and every $\sigma_i'$,
\[
  U_i(\sigma_i, \sigma_{-i}) \;\geq\; U_i(\sigma_i', \sigma_{-i}) - \varepsilon.
\]
\end{definition}

\subsection{Nash's Existence Theorem}

\begin{theorem}[Nash, 1950]\label{thm:nash}
Every finite normal-form game possesses at least one mixed-strategy Nash
equilibrium.
\end{theorem}

\begin{proof}[Proof sketch via Kakutani's fixed-point theorem]
Define the \emph{best-response correspondence}
$\br_i(\sigma_{-i}) = \argmax_{\sigma_i' \in \Delta(A_i)}
U_i(\sigma_i', \sigma_{-i})$.  Let $F : \Sigma \rightrightarrows \Sigma$
be the joint best-response map
$F(\sigma) = \prod_i \br_i(\sigma_{-i})$ where
$\Sigma = \prod_i \Delta(A_i)$.  The set $\Sigma$ is a non-empty,
compact, convex subset of Euclidean space.  Each $\br_i$ is
upper-hemicontinuous (by the Maximum Theorem) and convex-valued (since
$U_i$ is linear in $\sigma_i$).  Hence $F$ satisfies the hypotheses of
\emph{Kakutani's fixed-point theorem}, so there exists
$\sigma^* \in F(\sigma^*)$, which is by construction a Nash equilibrium.
\end{proof}

\subsection{Computational Complexity of Nash Equilibria}

\begin{theorem}[Daskalakis, Goldberg, Papadimitriou, 2009; Chen, Deng, 2006]
Computing a Nash equilibrium of a two-player normal-form game is
\emph{PPAD-complete}.
\end{theorem}

\begin{remark}
PPAD (Polynomial Parity Arguments on Directed graphs) is a complexity
class that captures the difficulty of problems whose existence is
guaranteed by a parity argument.  PPAD-completeness implies that, unless
PPAD$\,\subseteq\,$P (widely believed false), there is no polynomial-time
algorithm for finding an exact NE.  This motivates the search for
\emph{approximate} equilibria and equilibria of \emph{structured} game
classes.
\end{remark}

\subsection{Correlated Equilibrium}

\begin{definition}[Correlated equilibrium, Aumann 1974]
A distribution $\mu \in \Delta(\prod_i A_i)$ is a \emph{correlated
equilibrium} (CE) if, for every player~$i$ and every function
$\phi_i : A_i \to A_i$,
\[
  \sum_{a \in \prod_j A_j} \mu(a)\, u_i(a)
  \;\geq\;
  \sum_{a \in \prod_j A_j} \mu(a)\, u_i(\phi_i(a_i), a_{-i}).
\]
\end{definition}

The set of correlated equilibria is a \emph{convex polytope} defined by
linear inequalities, so a CE can be found in polynomial time via linear
programming.  Every NE induces a CE (take $\mu = \prod_i \sigma_i^*$),
but the converse is false; the CE polytope is generally larger.

\subsection{Coarse Correlated Equilibrium and No-Regret Learning}

\begin{definition}[Coarse correlated equilibrium]
A distribution $\mu$ is a \emph{coarse correlated equilibrium} (CCE) if,
for every player~$i$ and every action $a_i' \in A_i$,
\[
  \E_{a\sim\mu}[u_i(a)]
  \;\geq\;
  \E_{a\sim\mu}[u_i(a_i', a_{-i})].
\]
\end{definition}

The connection to learning is fundamental:

\begin{theorem}[No-regret $\Rightarrow$ CCE convergence]
\label{thm:no_regret_cce}
If every player~$i$ uses a no-external-regret learning algorithm (so that
each player's regret $R_i^T$ satisfies $R_i^T / T \to 0$), then the
empirical distribution of play converges to the set of coarse correlated
equilibria.
\end{theorem}

\begin{proof}[Proof sketch]
Let $\bar{\mu}^T = \frac{1}{T}\sum_{t=1}^{T} \delta_{a^t}$ be the
empirical distribution.  Player~$i$'s regret is
$R_i^T = \max_{a_i'}\sum_{t=1}^T [u_i(a_i', a_{-i}^t) - u_i(a^t)]$.
Dividing by $T$ and rearranging gives exactly the CCE condition with
slack $R_i^T/T \to 0$.
\end{proof}

\begin{remark}
If every player uses a no-\emph{internal}-regret (or no-\emph{swap}-regret)
algorithm, the empirical play converges to the set of \emph{correlated}
equilibria.  In two-player zero-sum games, no-regret learning yields
convergence to Nash equilibrium of the \emph{average} strategies (though
not the last iterate in general).
\end{remark}

%===================================================================
\section{Counterfactual Regret Minimisation (CFR)}\label{sec:cfr}
%===================================================================

\subsection{Extensive-Form Games}

\begin{definition}[Extensive-form game of imperfect information]
An extensive-form game is specified by:
\begin{itemize}
  \item A game tree with \emph{histories} $\cH$, terminal histories
    $\cZ \subset \cH$, and a player function $P : \cH\setminus\cZ \to
    N \cup \{c\}$ (where $c$ denotes chance).
  \item For each player~$i$, an \emph{information partition} $\cI_i$
    that groups histories into \emph{information sets}; at histories in
    the same information set, the player cannot distinguish the game
    states.
  \item Action sets $A(I)$ for each information set $I \in \cI_i$, with
    $A(h) = A(I)$ for all $h \in I$.
  \item Utility functions $u_i : \cZ \to \R$.
\end{itemize}
\end{definition}

A \emph{behavioural strategy} $\sigma_i$ assigns a distribution over
$A(I)$ to each information set $I \in \cI_i$.  By Kuhn's theorem, in
games of perfect recall, behavioural and mixed strategies are equivalent.

\paragraph{Notation.}  For a terminal history $z$ and strategy profile
$\sigma$, write $\pi^\sigma(z) = \prod_{h \sqsubset z} \sigma_{P(h)}(a_h)$
for the probability of reaching $z$, where $a_h$ is the action taken at
$h$ on the path to $z$.  Decompose:
\[
  \pi^\sigma(z) = \pi_i^\sigma(z)\cdot \pi_{-i}^\sigma(z)\cdot \pi_c(z),
\]
where $\pi_i^\sigma(z)$ is the product of player~$i$'s action
probabilities, $\pi_{-i}^\sigma$ is the opponents', and $\pi_c$ is
chance's.

\subsection{Counterfactual Value}

\begin{definition}[Counterfactual value]
The \emph{counterfactual value} of information set $I \in \cI_i$ under
profile $\sigma$ is
\[
  v_i^\sigma(I)
  = \sum_{z \in \cZ_I} \pi_{-i}^\sigma(z[I])\, \pi_c(z[I])\,
    \pi^\sigma(z \mid z[I])\, u_i(z),
\]
where $\cZ_I$ is the set of terminal histories passing through $I$,
$z[I]$ is the prefix of $z$ at $I$, and $\pi^\sigma(z \mid z[I])$ is the
probability of the continuation from $z[I]$ to $z$ under $\sigma$.
Intuitively, this is the expected value for player~$i$ at $I$, weighted
by the opponents' and chance's probability of reaching $I$, but
\emph{assuming player~$i$ plays to reach~$I$}.
\end{definition}

For a specific action $a \in A(I)$:
\[
  v_i^\sigma(I, a)
  = \sum_{z \in \cZ_{I,a}} \pi_{-i}^\sigma(z[I])\, \pi_c(z[I])\,
    \pi^\sigma(z \mid z[I], a)\, u_i(z).
\]

\subsection{Counterfactual Regret and Regret Matching}

\begin{definition}[Instantaneous counterfactual regret]
For action $a$ at information set $I$ at iteration~$t$:
\[
  r^t(I, a) = v_i^{\sigma^t}(I, a) - v_i^{\sigma^t}(I).
\]
The \emph{cumulative counterfactual regret} is
$R^T(I, a) = \sum_{t=1}^T r^t(I, a)$.
\end{definition}

\paragraph{Regret matching.}  At each iteration, CFR updates the strategy
at each information set via \emph{regret matching}:
\[
  \sigma^{T+1}(I, a) =
  \begin{cases}
    \dfrac{R^{T,+}(I, a)}{\displaystyle\sum_{a'} R^{T,+}(I, a')}
    & \text{if } \sum_{a'} R^{T,+}(I, a') > 0, \\[8pt]
    \dfrac{1}{|A(I)|} & \text{otherwise},
  \end{cases}
\]
where $R^{T,+}(I,a) = \max(R^T(I,a), 0)$.

\begin{algorithm}[t]
\caption{Vanilla CFR}
\label{alg:cfr}
\begin{algorithmic}[1]
\State Initialise $\sigma^1$ uniformly, $R^0(I,a) \gets 0$,
  $S^0(I,a) \gets 0$ for all $I, a$.
\For{$t = 1, 2, \ldots, T$}
  \For{each player $i$}
    \State Traverse game tree; compute $v_i^{\sigma^t}(I,a)$ for all
      $I \in \cI_i$, $a \in A(I)$.
    \State $R^t(I,a) \gets R^{t-1}(I,a) + r^t(I,a)$
    \State $S^t(I,a) \gets S^{t-1}(I,a) + \pi_i^{\sigma^t}(I)\,\sigma^t(I,a)$
      \Comment{Average strategy weights}
  \EndFor
  \State Update $\sigma^{t+1}$ via regret matching.
\EndFor
\State \Return $\bar{\sigma}^T(I,a) = S^T(I,a) / \sum_{a'} S^T(I,a')$
  \Comment{Average strategy}
\end{algorithmic}
\end{algorithm}

\subsection{Convergence Theorem}

\begin{theorem}[CFR convergence, Zinkevich et al.\ 2007]\label{thm:cfr}
In a two-player zero-sum extensive-form game of perfect recall, the
average strategy profile $\bar{\sigma}^T$ produced by CFR is a
$2\Delta_u \sqrt{|\cI| \cdot \max_I|A(I)|}\,/\,\sqrt{T}$\,-Nash
equilibrium, where $\Delta_u = \max_z u(z) - \min_z u(z)$ is the range
of utilities and $|\cI| = |\cI_1| + |\cI_2|$.
\end{theorem}

\begin{proof}[Proof sketch]
The key insight is the \emph{counterfactual regret decomposition}: the
overall regret of player~$i$ is bounded by the sum of positive
counterfactual regrets:
\[
  R_i^T \;\leq\; \sum_{I \in \cI_i} R^{T,+}_{\max}(I),
  \qquad
  R^{T,+}_{\max}(I) = \max_{a} R^{T,+}(I,a).
\]
Regret matching at each information set guarantees
$R^{T,+}_{\max}(I) \leq \Delta_u \sqrt{|A(I)| \cdot T}$ (this follows
from the regret bound of the multiplicative-weights / regret-matching
algorithm).  Therefore:
\[
  R_i^T
  \leq \sum_{I \in \cI_i} \Delta_u \sqrt{|A(I)|\, T}
  \leq \Delta_u \sqrt{T} \cdot |\cI_i| \cdot \sqrt{\max_I |A(I)|}.
\]
In a two-player zero-sum game, if both players have average regret at
most $\varepsilon$, then the average strategy profile is a
$2\varepsilon$-Nash equilibrium (since the game value is well-defined and
each player's average strategy is $\varepsilon$-close to a best
response).  The bound $\varepsilon = O(1/\sqrt{T})$ follows.
\end{proof}

\subsection{CFR+ and Variants}

\textbf{CFR+} (Tammelin et al., 2015; Bowling et al.) modifies the regret
accumulation rule to
\[
  R^t(I,a) \gets \max\bigl(R^{t-1}(I,a) + r^t(I,a),\; 0\bigr),
\]
i.e., cumulative regrets are clipped to non-negative values at every
step (``regret matching+'').  Additionally, the average strategy uses
linear weighting $S^t(I,a) \gets S^{t-1}(I,a) + t\cdot\pi_i^{\sigma^t}(I)\,\sigma^t(I,a)$, which upweights later (better) strategies.  Empirically, CFR+
converges orders of magnitude faster than vanilla CFR.

\textbf{Monte Carlo CFR (MCCFR)} uses sampling to avoid full tree
traversals.  Variants include \emph{outcome sampling} (single trajectory)
and \emph{external sampling} (sample opponent/chance, traverse all of
player~$i$'s actions).

\textbf{Deep CFR} (Brown et al., 2019) replaces the regret tables with
neural networks.  At each CFR iteration, advantage values are computed
via traversals and stored in a reservoir buffer; a neural network is
trained to predict cumulative advantages $\hat{R}^T(I,a)$.  A separate
network approximates the average strategy.

%===================================================================
\section{Libratus and Pluribus}\label{sec:libratus}
%===================================================================

\subsection{Libratus: Superhuman Heads-Up No-Limit Texas Hold'em}

Libratus (Brown and Sandholm, 2018) defeated top professional players in
heads-up no-limit Texas Hold'em.  The system has three modules:

\begin{enumerate}
  \item \textbf{Blueprint computation.}  An abstracted version of the game
    is solved using Monte Carlo CFR.  Card abstraction groups
    strategically similar hands; bet abstraction discretises the
    continuous betting actions.
  \item \textbf{Subgame solving.}  During play, when the opponent takes
    an action outside the abstraction, Libratus constructs and solves a
    \emph{subgame} in real time.
  \item \textbf{Self-improver.}  After each day of play, Libratus adds
    new bet sizes to the abstraction at branches the opponents
    frequently visited, then re-solves.
\end{enumerate}

\subsection{Safe Subgame Solving}

Naive subgame solving can increase exploitability.  The key theoretical
contribution is \emph{safe} (or reach-aware) subgame solving:

\begin{theorem}[Safe subgame solving, Brown and Sandholm 2017]
Let $\sigma$ be a Nash equilibrium of the full game and let $S$ be a
subgame rooted at a public state.  If we re-solve $S$ subject to the
constraint that the value achieved for each opponent reach distribution
is at least as high as under $\sigma$, then the resulting strategy is
part of a Nash equilibrium of the full game.
\end{theorem}

\begin{remark}
In practice, one does not have an NE of the full game; one has the
blueprint $\hat{\sigma}$.  Subgame solving refines the blueprint locally,
and the ``safe'' constraints prevent any increase in exploitability
relative to $\hat{\sigma}$.
\end{remark}

\subsection{Pluribus: Six-Player No-Limit Hold'em}

Pluribus (Brown and Sandholm, 2019) extended superhuman play to
six-player poker.  Key innovations:

\begin{itemize}
  \item \textbf{Depth-limited search.}  Instead of solving to the end of
    the game, Pluribus searches only a few moves ahead and uses a learned
    \emph{value function} (the blueprint strategy's expected value) at
    leaf nodes.
  \item \textbf{Linear CFR.}  Regrets are weighted linearly by iteration
    number: $R^T(I,a) = \sum_{t=1}^T t \cdot r^t(I,a)$, discounting
    early (poor) iterations.
  \item \textbf{Multi-player subgame solving.}  In multiplayer games,
    Nash equilibria are not unique and not interchangeable.  Pluribus
    handles this by assuming opponents play according to a small number
    of ``blueprint-like'' strategies and solving against this restricted
    opponent model.
\end{itemize}

%===================================================================
\section{Cicero: AI for Diplomacy}\label{sec:cicero}
%===================================================================

\subsection{The Challenge of Diplomacy}

Diplomacy is a seven-player board game of negotiation and alliance
formation.  Unlike poker, it requires:
\begin{itemize}
  \item \textbf{Natural language communication} --- players negotiate via
    free-form messages.
  \item \textbf{Cooperation and betrayal} --- alliances are non-binding.
  \item \textbf{Very large action spaces} --- each turn involves
    simultaneous orders for multiple units.
\end{itemize}
No prior AI had achieved human-level play.

\subsection{Cicero's Architecture (Meta FAIR, 2022)}

Cicero (FAIR/Meta, led by Noam Brown et al.) integrates:

\begin{enumerate}
  \item \textbf{Strategic reasoning engine.}  A planning module based on
    \emph{piKL} (policy-regularised iterative Kullback--Leibler
    divergence).  Given a prior policy $\pi_0$ (a supervised-learning
    model trained on human games), piKL solves:
    \[
      \sigma_i^* = \argmax_{\sigma_i} \left[
        U_i(\sigma_i, \sigma_{-i})
        - \lambda\, \KL(\sigma_i \| \pi_0)
      \right],
    \]
    iterating best responses while staying close to human-like play.
    This is crucial: pure equilibrium play would be alien to human
    partners and undermine cooperation.

  \item \textbf{Dialogue model.}  A large language model (fine-tuned
    BART/R2C2) that generates natural language messages \emph{conditioned
    on the intended actions}.

  \item \textbf{Intent model.}  Predicts other players' intended actions
    from their messages, enabling Cicero to ``read between the lines.''

  \item \textbf{Filtering pipeline.}  Generated messages are filtered
    for consistency with the strategic plan, honesty (relative to
    Cicero's actual intentions), and coherence.
\end{enumerate}

\subsection{Signaling Games Formalization}

Diplomacy can be partially formalised through the lens of \emph{signaling
games} (Crawford and Sobel, 1982):

\begin{definition}[Signaling game]
A signaling game consists of a \emph{sender} with private type
$\theta \in \Theta$, a \emph{message space} $M$, and a \emph{receiver}
who takes action $a \in A$ after observing the message.  A
\emph{perfect Bayesian equilibrium} specifies:
\begin{itemize}
  \item Sender strategy $\mu : \Theta \to \Delta(M)$,
  \item Receiver strategy $\alpha : M \to \Delta(A)$,
  \item Receiver beliefs $\beta : M \to \Delta(\Theta)$ consistent with
    Bayes' rule on the equilibrium path.
\end{itemize}
\end{definition}

In Diplomacy, each player is simultaneously a sender and receiver.  The
``type'' corresponds to the player's strategic intent (which units to move
where), the messages are natural language, and the actions are the orders
submitted.  Cicero's architecture approximates an equilibrium of this
signaling game: the intent model estimates $\beta$, the strategic engine
computes $\alpha$ (best-response actions given beliefs), and the dialogue
model implements $\mu$ (generating messages consistent with intent).

A key subtlety is \emph{cheap talk} (messages are costless): in classical
cheap talk games, only ``babbling'' equilibria survive under misaligned
incentives.  Diplomacy partly escapes this because players have
\emph{partially aligned} interests (they need allies) and the game is
repeated (reputation matters).

%===================================================================
\section{Mean-Field Games}\label{sec:mfg}
%===================================================================

\subsection{Motivation: The $N \to \infty$ Limit}

Many real-world multi-agent systems involve a very large number of
interacting agents: financial markets, autonomous vehicle fleets, content
recommendation ecosystems.  Solving for Nash equilibria in $N$-player
games is intractable for large $N$.  \emph{Mean-field game theory}
(independently introduced by Lasry and Lions, 2006--2007, and Huang,
Malham\'e, and Caines, 2006) studies the limiting regime $N \to \infty$,
where each agent interacts with the \emph{population distribution}
rather than with individual agents.

\subsection{The Mean-Field Game PDE System}

Consider a continuum of agents, each with state $x \in \R^d$, evolving
according to the controlled stochastic differential equation:
\[
  dX_t = b(X_t, \alpha_t, m_t)\,dt + \sigma\,dW_t,
\]
where $\alpha_t$ is the control, $m_t$ is the population distribution at
time~$t$, and $W_t$ is a standard Brownian motion.  Each agent minimises
a cost functional:
\[
  J(\alpha; m) = \E\left[
    \int_0^T L(X_t, \alpha_t, m_t)\,dt + G(X_T, m_T)
  \right].
\]

A \emph{mean-field Nash equilibrium} is a pair $(\alpha^*, m^*)$ such
that:
\begin{enumerate}
  \item $\alpha^*$ is optimal for the cost $J(\cdot\,; m^*)$ (taking
    $m^*$ as given),
  \item $m^*$ is the distribution of the state process under $\alpha^*$.
\end{enumerate}

This fixed-point condition yields a coupled PDE system:

\begin{theorem}[Mean-field game PDE system, Lasry and Lions 2007]
Under regularity conditions, a mean-field Nash equilibrium is
characterised by the system
\begin{align}
  -\partial_t u - \nu\Delta u + H(x, \nabla u, m) &= 0
    &\quad &\text{(Hamilton--Jacobi--Bellman, backward)},
    \label{eq:hjb} \\
  \partial_t m - \nu\Delta m
    - \mathrm{div}\!\left(m\,\nabla_p H(x, \nabla u, m)\right) &= 0
    &\quad &\text{(Fokker--Planck / Kolmogorov, forward)},
    \label{eq:fp}
\end{align}
with terminal condition $u(T,x) = G(x, m_T)$ and initial condition
$m(0, \cdot) = m_0$, where $H$ is the Hamiltonian of the optimal control
problem and $\nu = \sigma^2/2$.
\end{theorem}

The HJB equation~\eqref{eq:hjb} runs \emph{backward} in time from the
terminal cost, determining the value function and hence the optimal
policy.  The Fokker--Planck equation~\eqref{eq:fp} runs \emph{forward},
propagating the population density under the optimal policy.  The
coupling makes this a challenging forward--backward system.

\subsection{Existence and Uniqueness: Lasry--Lions Monotonicity}

\begin{definition}[Lasry--Lions monotonicity condition]
The running cost $L$ and terminal cost $G$ satisfy the \emph{Lasry--Lions
monotonicity condition} if, for all distributions $m_1, m_2$,
\[
  \int_{\R^d} \bigl[F(x, m_1) - F(x, m_2)\bigr]\,d(m_1 - m_2)(x)
  \;\geq\; 0,
\]
where $F$ denotes the coupling terms through which $m$ enters the costs.
\end{definition}

\begin{theorem}[Uniqueness under monotonicity]
If the Lasry--Lions monotonicity condition holds strictly, then the
mean-field game PDE system has at most one solution, and hence the
mean-field Nash equilibrium is unique.
\end{theorem}

Intuitively, monotonicity means that agents are ``repelled'' by
crowding---costs increase where the population concentrates.  This is
natural in congestion games but violated in models with herding or
network effects.

\subsection{Mean-Field Multi-Agent Reinforcement Learning}

In the ML setting, the MFG framework is combined with RL:

\begin{itemize}
  \item \textbf{Mean-Field MARL} (Yang et al., 2018): each agent's
    Q-function depends on its own action and the mean action of
    neighbours, reducing the joint action space from $|A|^N$ to
    $|A| \times |A|$.
  \item \textbf{MF-OMO} (Gemp et al., DeepMind): online mirror descent
    in the space of population distributions.
  \item \textbf{Neural MFG} (Carmona and Lauri\`ere, 2019): solve the
    coupled PDE system using neural networks (physics-informed or via
    direct policy optimisation).
\end{itemize}

%===================================================================
\section{Opponent Shaping}\label{sec:opponent_shaping}
%===================================================================

\subsection{Beyond Best Response: Shaping Opponents' Learning}

Standard game-theoretic algorithms assume opponents are fixed or
adversarial.  \emph{Opponent shaping} considers agents that explicitly
account for their influence on other agents' learning dynamics.

\subsection{LOLA: Learning with Opponent-Learning Awareness}

Foerster et al.\ (2018) introduced LOLA.  Consider two agents with
parameters $\theta_1, \theta_2$ and differentiable expected returns
$V_1(\theta_1, \theta_2)$, $V_2(\theta_1, \theta_2)$.  A naive learner
updates:
\[
  \theta_1 \gets \theta_1 + \alpha\,\nabla_{\theta_1} V_1(\theta_1, \theta_2).
\]

LOLA anticipates that player~2 will perform a gradient step, so
player~1 optimises $V_1$ \emph{after} player~2's anticipated update:

\begin{definition}[LOLA update]
Player~1's LOLA gradient is:
\[
  \nabla_{\theta_1}^{\text{LOLA}} V_1
  = \nabla_{\theta_1} V_1
  + \alpha_2\, \nabla_{\theta_1}\bigl(\nabla_{\theta_2} V_2\bigr)^\top
    \nabla_{\theta_2} V_1,
\]
where $\alpha_2$ is player~2's learning rate.  The additional term
$\nabla_{\theta_1}(\nabla_{\theta_2} V_2)^\top \nabla_{\theta_2} V_1$
is a second-order correction that ``steers'' player~2's learning in a
direction beneficial to player~1.
\end{definition}

\begin{remark}
The LOLA update requires computing a Hessian--vector product
$\nabla_{\theta_1}(\nabla_{\theta_2} V_2)$, which is tractable via
automatic differentiation (forward-mode or reverse-mode with
\texttt{jax.grad}).  Higher-order extensions (LOLA-$k$, SOS) account for
$k$ levels of recursive reasoning.
\end{remark}

\paragraph{Results.}
In the Iterated Prisoner's Dilemma, LOLA agents learn tit-for-tat--like
strategies and achieve mutual cooperation, whereas naive learners
converge to mutual defection.

\subsection{Model-Free Opponent Shaping (M-FOS)}

Lu et al.\ (2022) introduced M-FOS, which formulates opponent shaping as
a \emph{meta-game}: a ``meta-agent'' learns a policy over the space of
inner-loop learning algorithms.  The meta-agent's action at each step is
a parameter update (or a choice of learning rule) for the inner agent.

Key features of M-FOS:
\begin{itemize}
  \item Does not require differentiating through the opponent's update
    (model-free).
  \item Uses RL (PPO) to optimise the meta-agent's policy.
  \item Recovers LOLA-like behaviour in matrix games and scales to more
    complex settings.
  \item Can discover novel shaping strategies not captured by fixed-order
    Taylor approximations.
\end{itemize}

%===================================================================
\section{Connections to Modern Agentic AI}\label{sec:agentic}
%===================================================================

The recent proliferation of LLM-based agents creates new multi-agent
settings where game-theoretic reasoning is essential.

\subsection{Multi-Agent LLM Systems}

When multiple LLM-powered agents interact---e.g., a buyer agent
negotiating with a seller agent, or multiple coding agents working on a
shared codebase---the interaction is inherently game-theoretic:
\begin{itemize}
  \item Each agent has its own objective (its user's instructions).
  \item Agents may have private information (user preferences, budget
    constraints).
  \item Communication is in natural language, creating a signaling game.
\end{itemize}

Cicero's architecture provides a template: strategic reasoning (what to
do) is separated from communication (what to say), with an intent model
bridging the two.

\subsection{Mechanism Design for AI}

\emph{Mechanism design} (``reverse game theory'') asks: how should we
design the rules of interaction to achieve desirable outcomes?

\begin{itemize}
  \item \textbf{AI constitutions as mechanisms.}  A system prompt or
    RLHF reward model defines the ``rules'' under which an LLM agent
    operates.  Designing these to be incentive-compatible (so the agent
    has no reason to deviate) is a mechanism design problem.
  \item \textbf{Multi-agent coordination.}  When multiple agents share a
    resource (API rate limits, shared memory, a code repository), a
    mechanism (e.g., a priority queue, a voting protocol) determines
    allocation.  The VCG mechanism and its computational approximations
    are directly relevant.
  \item \textbf{Automated mechanism design.}  Conitzer and Sandholm's
    line of work on computationally designing optimal mechanisms can be
    applied to design interaction protocols for AI agents.
\end{itemize}

\subsection{Safety as Game Theory}

\begin{itemize}
  \item \textbf{Red-teaming as adversarial games.}  Testing AI systems
    against adversarial attacks (jailbreaks, prompt injections) is a
    zero-sum game between attacker and defender.
  \item \textbf{Principal--agent models.}  The alignment problem---ensuring
    an AI agent acts in its principal's interest---is a classical
    contract theory / principal--agent problem with moral hazard (hidden
    action) and adverse selection (hidden type).
  \item \textbf{Cooperative AI} (Dafoe et al., 2020).  Designing agents
    that can credibly commit to cooperative strategies, even in the
    absence of binding agreements, requires game-theoretic tools
    (commitment devices, reputation, repeated game strategies).
\end{itemize}

\subsection{Auction Theory for LLM APIs}

LLM inference is a scarce resource.  Multiple users compete for compute.
This is naturally modelled as an auction:
\begin{itemize}
  \item Users have private valuations for faster/better inference.
  \item The provider allocates compute (batch priority, model routing).
  \item Auction mechanisms (first-price, second-price, combinatorial)
    determine pricing and allocation.
\end{itemize}
The theoretical framework of Myerson (1981) optimal auction design
applies.  Recent work explores dynamic auctions for streaming compute,
connecting to the bandit literature.

%===================================================================
\section{Open Problems}\label{sec:open}
%===================================================================

\begin{openproblem}[Last-iterate convergence in general games]
No-regret algorithms converge in average to (coarse) correlated
equilibria in general games and to Nash equilibria in two-player
zero-sum games.  Under what conditions do \emph{last-iterate}
convergence guarantees hold, and at what rate?  Recent work (e.g.,
optimistic gradient descent/ascent) shows last-iterate convergence in
zero-sum games, but the general-sum and multi-player cases remain
largely open.
\end{openproblem}

\begin{openproblem}[Scalable equilibrium computation beyond two-player zero-sum]
CFR and its variants converge to NE in two-player zero-sum games.  In
general-sum or multiplayer games, NE need not be unique, may be
non-interchangeable, and computing any NE is PPAD-hard.  What
practically relevant relaxations (e.g., team equilibria, extensive-form
correlated equilibria, trembling-hand equilibria) are both tractable and
strategically meaningful?
\end{openproblem}

\begin{openproblem}[Learning in mean-field games without monotonicity]
Most theoretical guarantees for MFG rely on the Lasry--Lions monotonicity
condition.  Real-world systems (social networks, financial markets)
often exhibit herding, network effects, and strategic complementarities
that violate monotonicity.  Can we design provably convergent learning
algorithms for non-monotone MFGs?
\end{openproblem}

\begin{openproblem}[Game-theoretic foundations of multi-agent LLM interaction]
When multiple LLM agents interact via natural language, what is the
appropriate solution concept?  Classical equilibria assume common
knowledge of rationality, but LLM agents may have bounded rationality,
prompt-dependent biases, and opaque objectives.  Developing a formal
theory---possibly drawing on bounded rationality, level-$k$ thinking, or
cognitive hierarchy models---is a major open problem.
\end{openproblem}

\begin{openproblem}[Mechanism design under LLM manipulation]
LLM-powered agents can generate persuasive but potentially deceptive
text.  Classical mechanism design assumes agents can misreport their
\emph{type} but cannot alter the mechanism itself.  When agents can
manipulate the mechanism designer (e.g., via prompt injection, social
engineering), the revelation principle breaks down.  How do we design
robust mechanisms in this setting?
\end{openproblem}

\begin{openproblem}[Opponent shaping at scale]
LOLA and M-FOS have been demonstrated in small games.  Scaling opponent
shaping to large multi-agent environments (e.g., economy simulations,
multi-agent robotics) with many heterogeneous agents is an open
challenge.  Connections to meta-learning and the mean-field limit may
provide tractable approximations.
\end{openproblem}

\begin{openproblem}[Sample complexity of game solving with function approximation]
Deep CFR and similar methods use neural function approximation.  What
are the sample complexity guarantees?  Can we characterise the
approximation--regret trade-off when the counterfactual value function
class has limited capacity?
\end{openproblem}

%===================================================================
\section{Annotated Bibliography}\label{sec:biblio}
%===================================================================

\begin{enumerate}[label={[\arabic*]}, leftmargin=*, itemsep=6pt]

\item \textbf{Nash, J.\ (1950).}  \emph{Equilibrium points in $n$-person
  games.}  Proceedings of the National Academy of Sciences.\\
  The foundational paper proving existence of mixed-strategy equilibria in
  finite games via Kakutani's fixed-point theorem.

\item \textbf{Aumann, R.\ (1974).}  \emph{Subjectivity and correlation in
  randomized strategies.}  Journal of Mathematical Economics.\\
  Introduces correlated equilibrium, a solution concept strictly more
  general than Nash equilibrium and computable via linear programming.

\item \textbf{Zinkevich, M., Johanson, M., Bowling, M., and Piccione, C.\
  (2007).}  \emph{Regret minimization in games with incomplete
  information.}  NeurIPS.\\
  The original CFR paper.  Proves that regret matching at each
  information set yields $O(1/\sqrt{T})$ convergence to NE in
  two-player zero-sum extensive-form games.  Michael Bowling's group at
  the University of Alberta.

\item \textbf{Bowling, M., Burch, N., Johanson, M., and Tammelin, O.\
  (2015).}  \emph{Heads-up limit hold'em poker is solved.}  Science.\\
  Uses CFR+ to essentially solve heads-up limit Texas Hold'em
  ($\sim\!10^{14}$ information sets), producing a near-perfect
  strategy.

\item \textbf{Brown, N.\ and Sandholm, T.\ (2018).}  \emph{Superhuman AI
  for heads-up no-limit poker: Libratus beats top professionals.}
  Science.\\
  Describes the Libratus system: blueprint via MCCFR, safe subgame
  solving, and self-improvement.  A landmark in AI game-playing.

\item \textbf{Brown, N.\ and Sandholm, T.\ (2019).}  \emph{Superhuman AI
  for multiplayer poker.}  Science.\\
  Pluribus: extends superhuman poker play to six players using
  depth-limited search and linear CFR.

\item \textbf{Moravčík, M., Schmid, M., Burch, N., et al.\ (2017).}
  \emph{DeepStack: Expert-level artificial intelligence in heads-up
  no-limit poker.}  Science.\\
  DeepMind/U Alberta collaboration.  Uses continual re-solving with deep
  neural networks for value estimation.

\item \textbf{Brown, N., Lerer, A., Gross, S., and Sandholm, T.\ (2019).}
  \emph{Deep counterfactual regret minimization.}  ICML.\\
  Replaces tabular regret storage with neural networks, enabling CFR to
  scale to very large games.

\item \textbf{FAIR/Meta et al.\ (2022).}  \emph{Human-level play in the
  game of Diplomacy by combining language models with strategic
  reasoning.}  Science.\\
  The Cicero paper.  Combines a dialogue module (LLM) with a strategic
  planning module (piKL) and an intent prediction model, achieving
  human-level Diplomacy play.  Led by Noam Brown.

\item \textbf{Lasry, J.-M.\ and Lions, P.-L.\ (2007).}  \emph{Mean field
  games.}  Japanese Journal of Mathematics.\\
  Foundational paper on mean-field games.  Derives the coupled HJB /
  Fokker--Planck PDE system and proves existence and uniqueness under
  monotonicity conditions.

\item \textbf{Huang, M., Malham\'e, R., and Caines, P.\ (2006).}
  \emph{Large population stochastic dynamic games: closed-loop
  McKean--Vlasov systems and the Nash certainty equivalence principle.}
  Communications in Information and Systems.\\
  Independent co-discovery of mean-field game theory, with a focus on
  linear-quadratic models and the ``Nash certainty equivalence''
  principle.

\item \textbf{Foerster, J., Chen, R., Al-Shedivat, M., et al.\ (2018).}
  \emph{Learning with opponent-learning awareness.}  AAMAS.\\
  Introduces LOLA: agents that differentiate through the opponent's
  anticipated learning step.  Shows emergence of cooperation in the
  Iterated Prisoner's Dilemma.

\item \textbf{Lu, C., Willi, T., Letcher, A., and Foerster, J.\ (2022).}
  \emph{Model-free opponent shaping.}  ICML.\\
  M-FOS: uses meta-learning to shape opponents without requiring
  differentiable opponent models.

\item \textbf{Silver, D., Schrittwieser, J., Simonyan, K., et al.\
  (2017).}  \emph{Mastering the game of Go without human knowledge.}
  Nature.\\
  AlphaGo Zero / AlphaZero.  Self-play + MCTS + deep neural networks.
  While not directly game-theoretic in the equilibrium-computation
  sense, the self-play framework implicitly converges towards
  minimax-optimal play in two-player zero-sum games.

\item \textbf{Lanctot, M., Lockhart, E., Lespiau, J.-B., et al.\ (2019).}
  \emph{OpenSpiel: A framework for reinforcement learning in games.}
  arXiv:1908.09453.\\
  Open-source framework (DeepMind) implementing a wide range of games
  and algorithms (CFR, MCCFR, neural fictitious self-play, etc.).
  Essential infrastructure for game-theoretic AI research.

\item \textbf{Daskalakis, C., Goldberg, P., and Papadimitriou, C.\
  (2009).}  \emph{The complexity of computing a Nash equilibrium.}
  SIAM Journal on Computing.\\
  Proves that computing a Nash equilibrium in a two-player game is
  PPAD-complete, establishing a fundamental barrier to efficient
  computation.

\item \textbf{Yang, Y., Luo, R., Li, M., et al.\ (2018).}  \emph{Mean
  field multi-agent reinforcement learning.}  ICML.\\
  Applies the mean-field approximation to MARL, reducing the
  exponential joint action space to a tractable form.

\item \textbf{Dafoe, A., Hughes, E., Bachrach, Y., et al.\ (2020).}
  \emph{Open problems in cooperative AI.}  arXiv:2012.08630.\\
  A research agenda (DeepMind, with Thore Graepel among the authors)
  for building AI systems that can cooperate with humans and other
  agents.  Identifies key challenges at the intersection of game theory
  and AI safety.

\item \textbf{Conitzer, V.\ and Sandholm, T.\ (2002).}  \emph{Complexity
  of mechanism design.}  UAI.\\
  Foundational work on automated mechanism design: computationally
  designing optimal mechanisms for given objectives, connecting game
  theory to algorithm design.

\item \textbf{Carmona, R.\ and Lauri\`ere, M.\ (2019).}  \emph{Convergence
  analysis of machine learning algorithms for the numerical solution of
  mean field control and mean field games.}  Annals of Applied
  Probability.\\
  Provides theoretical guarantees for neural-network-based solvers of
  mean-field game PDEs.

\end{enumerate}

%===================================================================
\section*{Acknowledgments}
%===================================================================

These notes draw on the work of many research groups, including Meta FAIR
(Noam Brown, Adam Lerer), CMU (Tuomas Sandholm), the University of
Alberta (Michael Bowling, Neil Burch), DeepMind (David Silver, Thore
Graepel, Marc Lanctot, Ian Gemp), and the University of Oxford (Jakob
Foerster).  We are grateful to the broader game theory and multi-agent
learning communities for the open-source tools (OpenSpiel) and benchmark
environments that make this research possible.

\end{document}
